\documentclass[11pt]{article}
\usepackage[T1]{fontenc}
\usepackage[margin=1in]{geometry}
\usepackage{amsmath,amssymb,amsthm}
\usepackage[hidelinks]{hyperref}
\newtheorem{theorem}{Theorem}
\title{A single-exponential lower bound for Chromatic Number\\by vertex cover number}
\author{Atlas\thanks{The proof in this note was found using ChatGPT.}}
\date{}
\begin{document}
\maketitle

Given a simple graph and a positive integer \(\ell\), Chromatic Number
asks whether the graph admits a proper coloring with at most \(\ell\)
colors. Let \(k\) be the minimum size of a vertex cover and \(N\) the
binary input length.

\hypertarget{cover-w1}{}\begin{theorem}\label{cover-w1}
If \(\mathrm{FPT}\ne\mathrm{W[1]}\), then no deterministic algorithm
decides Chromatic Number in time \(2^{\mathcal O(k)}N^{\mathcal O(1)}\).
\end{theorem}

\begin{proof}
Suppose such an algorithm exists. We use it to decide Clique in
fixed-parameter time with respect to the requested clique size \(h\).
Clique is \(\mathrm{W[1]}\)-complete~\cite{df}. We may assume
\(2\le h\le n\), where \(n\) is the number of vertices in the input
graph. Make \(h\) copies of the vertex set, and join two vertices
from different copies precisely when their originals are adjacent.
A clique choosing one vertex per copy exists exactly when the input
graph has an \(h\)-clique.

\paragraph{Encoding vertex choices.}
Put \(r=\lceil\log_2 n\rceil\), \(q=2^h\), \(b=4q^2\), and
\(L=b\lceil r/b\rceil\). Pad each copy to \(2^L\) vertices with
isolated vertices and assign its vertices distinct \(L\)-bit names.
This preserves the answer since \(h\ge2\). Moreover, \(L<r+b\), so
\(2^L<2^{b+1}n\). Choosing one name from each copy amounts to
choosing a column in \(\{0,1\}^h\) at each of the \(L\) bit positions.

Partition the positions into blocks of \(b\) consecutive positions.
In each block, guess the number of occurrences of each of the \(q\)
possible columns. These nonnegative counts must sum to \(b\).
Since \(q\log_2(b+1)\le b\), there are at most
\((b+1)^{qL/b}\le2^L\) guesses over all blocks, even if we enumerate
all vectors of counts and discard invalid sums. We construct one
instance of Chromatic Number per valid guess and accept if any answers
yes. For a fixed guess and a block \(H\), prepare \(b\) labels
indexed by \(1,\ldots,b\), each carrying one \(h\)-bit column, with the
guessed multiplicities. Permuting these labels over the positions
realizes precisely the sequences of columns with those counts.

\paragraph{Encoding permutations by colors.}
Our graph contains a clique with \(2L\) vertices, and we set the color
budget to \(2L\). Thus its vertices name all available colors. To allow
a vertex a specified set of colors, make it adjacent to the clique
vertices for all other colors. For every block \(H\), reserve two sets
of colors \(A_H=\{a_{H,1},\ldots,a_{H,b}\}\) and
\(D_H=\{d_{H,1},\ldots,d_{H,b}\}\). All these sets are disjoint.
For each position \(j\in H\) introduce vertices \(x_j,y_j\) with
allowed sets \(A_H,D_H\), respectively. Within \(H\), make the
\(x_j\) pairwise adjacent and make the \(y_j\) pairwise adjacent.
Each set of vertices therefore uses a permutation of the \(b\) colors.

We enforce the same permutation in both sets. For every \(j\in H\)
and every ordered pair \(s\ne t\) of indices, give a new vertex the
allowed colors \(\{a_{H,s},d_{H,t}\}\) and join it to \(x_j,y_j\).
Both colors are blocked exactly when \(x_j,y_j\) choose indices
\(s,t\), respectively. Once all these vertices can be colored, there
is a single permutation \(\pi_H\) from the positions of \(H\) to the
label indices, realized by both sets. Conversely, any such permutation
allows every one of these vertices to be colored.

\paragraph{Forbidding pairs of nonadjacent vertices.}
For each nonedge between a vertex with name \(u\) in copy \(i\) and a
vertex with name \(v\) in copy \(i'>i\), we want to forbid the
simultaneous choice of \(u\) and \(v\). For each block \(H\), let
\(R_H\) be the positions in \(H\) where \(u\) has a one, and
\(T_H\) the label indices whose column has a one in coordinate \(i\).
Define \(R'_H,T'_H\) in the same way for \(v,i'\).
The encoded names agree with these two names on the whole block
exactly when
\[
   \pi_H(R_H)=T_H \qquad\text{and}\qquad \pi_H(R'_H)=T'_H .
\]
If either equality of cardinalities fails in any block, the forbidden
pair cannot be selected with this guess, and we skip the nonedge.

Otherwise introduce one vertex \(z\) for this nonedge. Over all
blocks \(H\), its allowed colors are the \(a_{H,t}\) for \(t\in T_H\)
and the \(d_{H,t}\) for \(t\in T'_H\). Its neighbors outside the
palette clique are the \(x_j\) for \(j\in R_H\) and the \(y_j\) for
\(j\in R'_H\) over all blocks. Then \(z\) has no available color
exactly when \(T_H\subseteq\pi_H(R_H)\) and
\(T'_H\subseteq\pi_H(R'_H)\) in every block. Since the cardinalities
match, these inclusions are exactly the two required equalities.
Thus \(z\) has no available color exactly when both forbidden
vertices are selected. This also covers empty sets, since the
corresponding inclusion then holds for every permutation.

All vertices added to test permutations or nonedges are mutually
nonadjacent, so their colors can be chosen independently.
A clique choosing one vertex per copy yields a valid guess and
permutations that admit all test colors. Conversely, an accepting
coloring determines one name per copy and all selected vertices are
pairwise adjacent. Hence the original Clique instance is positive
exactly when some guess admits a proper coloring.

\paragraph{Size and running time.}
The palette clique and all \(x_j,y_j\) form a vertex cover of size
\(4L\). Each query has at most \(Lb^2\) permutation tests and
\(\binom{h}{2}2^{2L}\) nonedge tests. Each test has at most \(4L\)
neighbors. Since \(L<\lceil\log_2 n\rceil+b\) and \(b\) depends
only on \(h\), all queries can be generated deterministically in
\(f(h)n^{\mathcal O(1)}\) time for a computable \(f\), including
the explicit binary encoding of all edges. In particular \(N\) has
the same bound. For each query \(k\le4L\), so for any fixed constant
\(c>0\),
\[
  2^{ck}\le 2^{4c(b+1)}n^{4c}.
\]
There are fewer than \(2^{b+1}n\) queries. Running the assumed
algorithm on all of them gives a deterministic fixed-parameter
algorithm for Clique, a contradiction.
\end{proof}

\begin{thebibliography}{1}
\bibitem{df} R.~G. Downey and M.~R. Fellows.
Fixed-parameter tractability and completeness II: On completeness for W[1].
\emph{Theoretical Computer Science} 141(1--2), 109--131, 1995.
\href{https://doi.org/10.1016/0304-3975(94)00097-3}{doi:10.1016/0304-3975(94)00097-3}.
\end{thebibliography}
\end{document}
